% =====================================================================
%  S1 (Supplementary Notes): the annotation and review protocol suite.
%  Fragment \input by supplement.tex under the S1 part. No preamble here.
% =====================================================================
\section{The annotation and review protocol suite}
\label{sec:supp-protocol-suite}

The annotation workflow described in the main paper has been
consolidated into a versioned, openly published protocol suite,
available on GitHub in the project
repository.\footnote{\url{https://github.com/ForomePlatform/genetic-evidence-model/tree/master/protocols}}
The suite separates the rules of
annotation from their execution. A mode-agnostic annotation protocol
(\texttt{PROTOCOL.md}) states what a well-formed annotation requires;
two mode-specific protocols specialise it for autonomous drafting
(\texttt{PROTOCOL\_AUTONOMOUS.md}) and for interactive,
curator-in-the-loop work; and a separate review protocol
(\texttt{REVIEW\_PROTOCOL.md}, with an interactive variant) governs the
assertion-by-assertion review of drafted annotations. A
labelling-examples document records worked cases for the recurring
judgement calls. The protocols are packaged as two executable skills
(defined in \cref{sec:supp-skills}), one for annotation and one for
review, so that the same rules can be applied by a human curator or by
an AI assistant working under curator supervision.

Two points of provenance are recorded. First, the pilot corpus
described in the main paper was annotated under an earlier procedure
(protocol v0); the suite described here is the later formalisation, and
the pilot annotations remain valid under it. Second, the AI-drafted
annotations were drafted with Claude Opus 4.7 and reviewed with Claude
Opus 4.8 under the review skill by an author whose expertise is in
genomic data curation and the evidence model rather than in human
or population genetics, with every verdict adjudicated by the curator; the
curator, not the model, is the authority for the recorded annotations.
